\section{Calculer une expresssion avec ou sans parenthèses}

\subsection{Cas d'une expression sans parenthèses}

\begin{convention}[\hspace{5pt} \it{$\times$ et $\div$ prioritaires sur + et -}]
    Dans une expression sans parenthèses, les \bb{multiplications} et les
    \bb{divisions} doivent être effectuées \answer{avant} les \bb{additions} et les
    \bb{soustractions}. \\

    On dit que les multiplications et les divisions sont
    \answer{prioritaires} sur l'addition et la soustraction.
\end{convention}

\begin{exe}
    Calculer $3+2\times 4$ en \bb{soulignant} l'opération réalisée à chaque étape.
    \begin{align*}
        3+2\times 4 &= \answer{3+\underline{2 \times 4}}  \\
        &=\answer{\underline{3+8}}\\
        &=\answer{11}
    \end{align*}
\end{exe}


\begin{convention}[\hspace{5pt} \it{Seulement des + et des -}]
    Dans une expression sans parenthèses qui ne contient \bb{que des additions et
    des soustractions}, on effectue les calculs \answer{de gauche à droite} (on dit aussi
    \it{dans le sens de lecture}).
\end{convention}

\begin{exe}
    Calculer $10-3+6$ en \bb{soulignant} l'opération réalisée à chaque étape.
    \begin{align*}
        10-3+6 &= \answer{\underline{10-3}+6}  \\
        &=\answer{\underline{7+6}}\\
        &=\answer{13}
    \end{align*}
\end{exe}


\begin{convention}[\hspace{5pt} \it{Seulement des $\times$ et des $\div$}]
    Dans une expression sans parenthèses qui ne contient \bb{que des multiplications
    et des divisions} on effectue les calculs \answer{de gauche à droite} (on dit aussi
    \it{dans le sens de lecture}).
\end{convention}

\begin{exe}
    Calculer $30\div 5 \times 2$ en \bb{soulignant} l'opération réalisée à chaque étape.
    \begin{align*}
        30\div 5 \times 2 &= \answer{\underline{30\div 5} \times 2}  \\
        &=\answer{\underline{6\times 2}}\\
        &=\answer{12}
    \end{align*}
\end{exe}


\begin{prop}[\it{Commutativité de l'addition et de la multiplication}]
    Si dans une expression sans parenthèses il n'y a \bb{que des additions} ou
    \bb{que des multiplications} alors on peut effectuer les calculs dans
    \bb{l'ordre que l'on veut}. 
    \newline
    \newline
    On dit que l'addition et la multiplication sont \bb{commutatives.}
\end{prop}

\begin{exe}
    Pour les additions, on peut choisir l'ordre que l'on veut : 

    \begin{minipage}{0.48\textwidth}
        $$2+3+5 = \underline{2+3}+5 = \underline{5} + 5 = 10$$
    \end{minipage}
    \hfill
    \begin{minipage}{0.48\textwidth}
        $$2+3+5 = 2+ \underline{3+5}= 2 + \underline{8} = 10$$
    \end{minipage}
    \newline
    \newline
    \newline
    De même, pour les multiplications :

    \begin{minipage}{0.48\textwidth}
        $$5\times37\times2 = \underline{5 \times 37} \times 2 = \underline{185} \times 2 = 370$$
    \end{minipage}
    \hfill
    \begin{minipage}{0.48\textwidth}
        $$5\times37\times2 = \underline{5} \times 37 \times \underline{2} = \underline{10}
        \times 37 = 370$$
    \end{minipage}
\end{exe}


\begin{aretenir}
    Il est important de \bb{toujours souligner/surligner} à chaque étape le 
    calcul prioritaire : seul ce qui est souligné/surligné est calculé à
    l'étape d’après, le reste n'est pas utilisé et doit donc être recopié
    \bb{au même endroit}.
\end{aretenir}

\subsection{Cas d'une expression avec parenthèses}
\subsubsection{Parenthèses simples}

\begin{convention}
    On commence par effectuer les calculs \bb{contenus dans les parenthèses.} 
    Même si elles sont à la fin du calcul.
\end{convention}

\begin{exe}
    \begin{align*}
        3 \times (2 + 4) &= \answer{3 \times (\underline{2+4})}\\
        &= \answer{\underline{3 \times (6)}}\\
        &= \answer{18}
    \end{align*}
\end{exe}
\subsubsection{Parenthèses doubles}
\begin{convention}
    On commmence par les calculs contenus dans les
    \answer{\bb{parenthèses les plus intérieures.}} 
\end{convention}
\begin{exe}
    \begin{align*}
        3 \times (2 \times (2+1)) &= \answer{3 \times (2 \times (\underline{2+1}))}\\
        &= \answer{3 \times (\underline{2\times3})}\\
        &= \answer{\underline{3\times6}}\\
        &= \answer{18}
    \end{align*}
\end{exe}

\subsubsection{Cas particulier : la fraction}
\begin{remarque}
    Un quotient peut se noter sous forme linéaire avec le signe $\div$ ou sous
    forme fractionnaire. Par exemple: 
    $$\dfrac{30}{5}=30\div5=6$$
    
    On peut alors calculer les fractions en suivant deux méthodes. Prenons
    l'exemple suivant :
    $$A = \dfrac{12\times2-21}{4+2}$$
    \begin{minipage}{0.49\textwidth}
    \begin{center}
        Ecriture fractionnaire
    \end{center}
        \begin{align*}
            A&=\answer{\dfrac{\hl{12\times2}-21}{\hl{4+2}}}\\
            A&=\answer{\dfrac{\hl{24-21}}{6}}\\
            A&=\answer{\dfrac{3}{6}}\\
            A&=\answer{0.5}
        \end{align*}
    \end{minipage}
    \vline
    \begin{minipage}{0.49\textwidth}
    \begin{center}
        Ecriture linéaire
    \end{center}
   \begin{align*}
       A&=\answer{(\hl{12\times2}-21)\div(\hl{4+2})}\\
        A&=\answer{(\hl{24-21})\div(6)}\\
        A&=\answer{(3)\div(6)}\\
        A&=\answer{0.5}
    \end{align*}
    \end{minipage}
\end{remarque}

\begin{attention}
    Les parenthèses sont \bb{obligatoires} lorsqu'on passe de l'écriture
    fractionnaire à l'écriture linéaire !
\end{attention}

\begin{align*}
    $\frac{56}{4}$  
\end{align*}
